\documentclass[a4paper,11pt]{ltjsarticle}
\usepackage{luatexja}
\usepackage{luatexja-fontspec}
\usepackage{amsmath,amssymb,amsthm}
\usepackage{mathtools}
\usepackage{booktabs}
\usepackage{longtable}
\usepackage{array}
\usepackage{hyperref}
\usepackage{graphicx}
\usepackage{float}
\usepackage{geometry}
\geometry{margin=25mm}

\hypersetup{
  colorlinks=true,
  linkcolor=blue,
  urlcolor=blue,
  citecolor=blue
}

\theoremstyle{definition}
\newtheorem{theorem}{定理}[section]
\newtheorem{proposition}[theorem]{命題}
\newtheorem{lemma}[theorem]{補題}
\newtheorem{corollary}[theorem]{系}
\newtheorem{definition}[theorem]{定義}
\newtheorem{remark}[theorem]{注意}
\newtheorem{observation}[theorem]{観察}

\setlength{\parindent}{1em}
\setlength{\parskip}{0.5em}

\providecommand{\tightlist}{\setlength{\itemsep}{0pt}\setlength{\parskip}{0pt}}
\providecommand{\real}[1]{#1}
\begin{document}

\section{\texorpdfstring{4次元球の離散構造：単位立方体充填と漸近的体積不足
\(\Delta(R)\)（論文3：離散性の数学的基盤）}{4次元球の離散構造：単位立方体充填と漸近的体積不足 \textbackslash Delta(R)（論文3：離散性の数学的基盤）}}\label{ux6b21ux5143ux7403ux306eux96e2ux6563ux69cbux9020ux5358ux4f4dux7acbux65b9ux4f53ux5145ux586bux3068ux6f38ux8fd1ux7684ux4f53ux7a4dux4e0dux8db3-deltarux8ad6ux65873ux96e2ux6563ux6027ux306eux6570ux5b66ux7684ux57faux76e4}

\textbf{著者}：木原 範昭（WF System Co., Ltd.）
\textbf{ORCID}：\href{https://orcid.org/0009-0004-6753-4020}{0009-0004-6753-4020}
\textbf{日付}：2026年4月
\textbf{DOI}：\href{https://doi.org/10.5281/zenodo.19837592}{10.5281/zenodo.19837592}

\begin{center}\rule{0.5\linewidth}{0.5pt}\end{center}

\subsection{要旨}\label{ux8981ux65e8}

半径 \(R = 2k+1\) の4次元球 \(B(R)\)
への、整数格子点を中心とする単位立方体の充填を考察する。\(B(R)\)
に完全包含される単位立方体の個数 \(N(k)\) を \(k \le 60\)
まで厳密に計算する。体積不足 \(\Delta(R) := V_4(R) - N(k)\) が、補助領域
\(\Omega_R = \{x \in \mathbb{R}^4 : \sum (|x_i|+1/2)^2 \le R^2\}\)
に対する包除原理から、\(R \to \infty\) で漸近展開
\[\Delta(R) = (16\pi/3)\,R^3 - 6\pi\,R^2 + O(R)\]
を持つことを示す。主要係数は
\[c := \lim_{k\to\infty} \Delta(R)/(2\pi^2 R^3) = 8/(3\pi) \approx 0.84883\]
となり、多項式最小二乗フィットによる数値検証で 0.024\%
以内の精度で一致する。充填数 \(N(k)\) は、変数変換 \(p_i = 2|x_i|+1\)
により4個の正の奇数平方和への制約に帰着し、Lagrange--Jacobi
の四平方理論と接続する。

\begin{center}\rule{0.5\linewidth}{0.5pt}\end{center}

\subsection{§1. 序論}\label{ux5e8fux8ad6}

古典的な Gauss 円問題は、半径 \(R\) の円盤内の整数点
\((m_1, m_2) \in \mathbb{Z}^2\) の個数 \(N_2(R)\) を問う。主要項は面積
\(\pi R^2\) で、誤差項のサイズ --- \(N_2(R) = \pi R^2 + O(R^\theta)\)
における最適指数 \(\theta\) ---
が長年の研究対象である。4次元類似は、Lagrange の四平方定理と Jacobi
による表現数 \(r_4(N)\)
の閉形式により、整数論的に完全に解かれた問題である。

本論文では、これと関連するが異なる問いを扱う。\(B(R)\)
内の整数点を数える代わりに、整数格子点 \(c \in \mathbb{Z}^4\)
を中心とする\textbf{単位立方体} --- 積
\(\prod_i [c_i - 1/2, c_i + 1/2]\) --- のうち、\(B(R)\)
に\textbf{完全に含まれる}ものの個数を数える。包含条件は不等式
\(\sum_i (|c_i| + 1/2)^2 \le R^2\) となり、これは整数点計数の単純な条件
\(\|c\|_2 \le R\) より厳密に厳しい。

動機は幾何学的である：単位立方体充填は、球 \(B(R)\) のどれだけが辺長
\(1\) の正則離散構造で充填可能かを測る。体積不足
\[\Delta(R) := V_4(R) - N(k), \qquad V_4(R) = (\pi^2/2)\,R^4\]
は、そのような充填では到達できない \(B(R)\)
の部分を定量化する。本論文の主要結果は、\(\Delta(R)\)
が精密に決定された漸近係数を持つことである：
\[\Delta(R) \sim \frac{16\pi}{3}\,R^3, \qquad c := \lim_{k\to\infty} \frac{\Delta(R)}{2\pi^2 R^3} = \frac{8}{3\pi} \approx 0.84883.\]

定数 \(c\) は補助領域
\(\Omega_R = \{x \in \mathbb{R}^4 : \sum (|x_i|+1/2)^2 \le R^2\}\)
に対する包除原理の計算により解析的に導出され、\(\Omega_R\)
内の格子点数は \(N(k)\) と高次誤差で一致する。\(N(k)\) の
\(k = 0, 1, \ldots, 60\)
までの数値計算は、多項式最小二乗フィットにより解析結果を 0.024\%
以内の精度で確認する。

論文構成。第2節で記号と問題設定を固定する。第3節で包含不等式とその整数定式化を導出する。第4節で
\(N(k)\) の基本的組合せ論的性質、特に角の立方体 16 個が \(\partial B\)
に乗る恒等式と漸近充填密度の収束を記述する。第5節で \(\Delta(R)\)
の漸近解析と \(c = 8/(3\pi)\) の導出を行う。第6節で Lagrange--Jacobi
接続を展開する。第7節で数値検証を提示する。第8節で結論を述べる。

本論文は純粋に数学的内容である。物理的解釈は提供しない。本結果の帰結（4+1次元ブラックホール熱力学との接続可能性を含む）は別所で扱う。

\begin{center}\rule{0.5\linewidth}{0.5pt}\end{center}

\subsection{§2. 問題設定}\label{ux554fux984cux8a2dux5b9a}

\subsubsection{§2.1 4次元球}\label{ux6b21ux5143ux7403}

\(R > 0\) を実数とする。原点中心の半径 \(R\) の4次元球を
\[B(R) := \{x = (x_1, x_2, x_3, x_4) \in \mathbb{R}^4 : \|x\|_2 \le R\}\]
とする（\(\|x\|_2 = (x_1^2 + x_2^2 + x_3^2 + x_4^2)^{1/2}\)）。

体積と境界 \(\partial B(R) \simeq S^3(R)\) の3次元体積：
\[V_4(R) = \frac{\pi^2}{2} R^4, \qquad S_3(R) = 2\pi^2 R^3.\]

これらは \(B(R)\)
に関連する整数計数の漸近挙動を扱う際の自然な正規化として繰り返し現れる。

\subsubsection{§2.2
単位立方体充填問題}\label{ux5358ux4f4dux7acbux65b9ux4f53ux5145ux586bux554fux984c}

中心 \(c \in \mathbb{R}^4\) の \(\mathbb{R}^4\)
中の\textbf{単位立方体}を
\[Q(c) := \{x \in \mathbb{R}^4 : |x_i - c_i| \le \tfrac{1}{2}, \;\; i = 1, 2, 3, 4\}\]
とする。\(Q(c)\) は辺長 \(1\)、4次元体積 \(1\)。\(2^4 = 16\) 個の頂点は
\(c + (\pm 1/2, \pm 1/2, \pm 1/2, \pm 1/2)\)。

中心が整数格子点 \(c \in \mathbb{Z}^4\) にある立方体に注目する。問い：

\textbf{(a)} どの \(c \in \mathbb{Z}^4\) で単位立方体 \(Q(c)\) が球
\(B(R)\) に完全包含されるか？

\textbf{(b)} \(R\) の関数として、そのような立方体は何個あるか？

\begin{enumerate}
\def\labelenumi{(\alph{enumi})}
\setcounter{enumi}{1}
\tightlist
\item
  の答えを \(N_4(R)\) と書く：
  \[N_4(R) := \#\{c \in \mathbb{Z}^4 : Q(c) \subseteq B(R)\}.\]
\end{enumerate}

これは \(B(R)\) にはみ出さずに収まる、整数中心単位立方体の個数。

\subsubsection{§2.3
奇数直径への制限}\label{ux5947ux6570ux76f4ux5f84ux3078ux306eux5236ux9650}

組合せ論的便宜のため、\(R\) を
\[R_k := 2k + 1, \qquad k \in \mathbb{Z}_{\ge 0}\]
の値に制限する。動機：

\begin{enumerate}
\def\labelenumi{\arabic{enumi}.}
\item
  \textbf{対称性}：\(Q(c) \subseteq B(R)\) は \(c_i \mapsto -c_i\)
  で不変。\(R = 2k+1\) なら、中心 \(c = 0\) は全 \(k \ge 0\)
  で含まれ（\(B(1)\)
  に原点周りの単位立方体が収まる）、格子点は曖昧さなく対称多重項を成す。
\item
  \textbf{整数演算}：充填不等式の両辺に \(4\)
  を掛けると全量が整数になり、厳密計算可能。
\end{enumerate}

本論文では \(R = R_k = 2k + 1\) とし、\(N(k) := N_4(R_k)\)
を充填数と呼ぶ。

この制限は本質的でない：全整数 \(R\) や連続 \(R\)
への拡張は容易な変更で可能。漸近結果は変更なしで成立し、変わるのはデータ点の密度のみ。

\begin{center}\rule{0.5\linewidth}{0.5pt}\end{center}

\subsection{§3.
充填条件の定式化}\label{ux5145ux586bux6761ux4ef6ux306eux5b9aux5f0fux5316}

\subsubsection{§3.1 包含不等式}\label{ux5305ux542bux4e0dux7b49ux5f0f}

中心 \(c \in \mathbb{R}^4\) の立方体 \(Q(c)\) が \(B(R)\)
に含まれるのは、その16個の頂点全てが \(B(R)\)
にあるとき、かつそのときに限る。16頂点のうち原点から最も遠いのは
\[P_{\max}(c) = (|c_1| + \tfrac{1}{2}, |c_2| + \tfrac{1}{2}, |c_3| + \tfrac{1}{2}, |c_4| + \tfrac{1}{2})\]
である（各座標の絶対値は \(c_i\) と同符号の \(1/2\)
を加えると最大化される。\(c_i = 0\) ならどちらの符号でもよい）。

包含条件 \(Q(c) \subseteq B(R)\) は \(\|P_{\max}(c)\|_2 \le R\) と同値：

\textbf{補題 3.1}。\emph{\(c \in \mathbb{Z}^4\), \(R > 0\)
とする。\(Q(c) \subseteq B(R)\) が成り立つのは}
\[\boxed{\;\sum_{i=1}^{4} (|c_i| + \tfrac{1}{2})^2 \le R^2\;}\]
\emph{が成り立つとき、かつそのときに限る。}

\textbf{証明}。\(Q(c)\) の頂点は
\(c + (\epsilon_1/2, \epsilon_2/2, \epsilon_3/2, \epsilon_4/2)\)（\(\epsilon_i \in \{-1, +1\}\)）。頂点のノルム平方は
\[\sum_i (c_i + \epsilon_i/2)^2 = \sum_i (c_i^2 + \epsilon_i c_i + 1/4)\]
で、\(\epsilon\) にわたる最大は
\(\epsilon_i = \mathrm{sgn}(c_i)\)（\(c_i \neq 0\) の場合、\(c_i = 0\)
ならどちらでも）で達成され、最大値は
\(\sum_i (|c_i| + 1/2)^2\)。全頂点が \(B(R)\) にある条件は、この最大値が
\(R^2\) を超えないこと。\(\square\)

\subsubsection{§3.2 整数定式化}\label{ux6574ux6570ux5b9aux5f0fux5316}

\(p_i := 2|c_i| + 1\) と置くと補題 3.1 の不等式は
\[\sum_{i=1}^{4} \left(\frac{p_i}{2}\right)^2 \le R^2 \qquad \text{すなわち} \qquad \sum_{i=1}^{4} p_i^2 \le (2R)^2.\]
\(R = 2k+1\) に対し \[\sum_{i=1}^{4} p_i^2 \le (4k+2)^2\] となる。各
\(p_i \in \{1, 3, 5, \ldots\}\) は正の奇整数。

この整数定式化には2つの帰結がある。第一に、充填問題が古典的平方和数論の枠内に位置づけられる（§6
で活用）。第二に、任意精度整数演算で厳密計算可能となり、浮動小数点誤差なしに進められる。

\subsubsection{\texorpdfstring{§3.3 充填数
\(N(k)\)}{§3.3 充填数 N(k)}}\label{ux5145ux586bux6570-nk}

補題 3.1 と \(c \in \mathbb{Z}^4\) から
\((p_i) \in \{1, 3, 5, \ldots\}^4\) への変換を組み合わせて
\[N(k) = \sum_{\substack{c \in \mathbb{Z}^4 \\ \sum (|c_i| + 1/2)^2 \le R_k^2}} 1.\]

奇正整数 \(p_i\) で表現すると：
\[N(k) = \sum_{\substack{(p_1, \ldots, p_4) \in \{1,3,5,\ldots\}^4 \\ \sum p_i^2 \le (4k+2)^2}} 2^{\#\{i : p_i > 1\}}.\]

因子 \(2^{\#\{p_i > 1\}}\) は
\(c_i = \pm (p_i - 1)/2\)（\(p_i > 1\)、\(c_i \neq 0\)）の符号自由度を反映する。\(p_i = 1\)
なら \(c_i = 0\) で符号自由度なし。

\subsubsection{§3.4
計算上の考察}\label{ux8a08ux7b97ux4e0aux306eux8003ux5bdf}

\(\mathbb{Z}^4\) の \(B(R)\) 内の直接列挙は \(O(R^4)\)
コスト。\(R = 2k+1\)
で本論文の範囲（\(k \le 60\)、\(R \le 121\)）では計算可能だが効率的でない。

より効率的な列挙：非負・非増加順序組
\(y_1 \ge y_2 \ge y_3 \ge y_4 \ge 0\)
に制限し、符号と順列の倍数を解析的に扱い、立方体の超八面体対称群
\(|B_4| = 384\) の倍率でコストを削減する：

\[N(k) = \sum_{\substack{y_1 \ge y_2 \ge y_3 \ge y_4 \ge 0 \\ \sum (y_i + 1/2)^2 \le R_k^2}} \mu_{\mathrm{sign}}(y) \cdot \mu_{\mathrm{perm}}(y),\]

ここで
\[\mu_{\mathrm{sign}}(y) = 2^{\#\{i : y_i > 0\}}, \qquad \mu_{\mathrm{perm}}(y) = \frac{4!}{\prod_v |\{i : y_i = v\}|!}.\]

この列挙は付属コードで実装され、\(k \le 60\) の全 \(N(k)\)
を市販ハードウェアで1分未満で正確に与える。

\begin{center}\rule{0.5\linewidth}{0.5pt}\end{center}

\subsection{\texorpdfstring{§4. 数列
\(N(k)\)：組合せ論的性質}{§4. 数列 N(k)：組合せ論的性質}}\label{ux6570ux5217-nkux7d44ux5408ux305bux8ad6ux7684ux6027ux8cea}

\subsubsection{§4.1 初期値}\label{ux521dux671fux5024}

直接計算により：

{\def\LTcaptype{none} % do not increment counter
\begin{longtable}[]{@{}rrrr@{}}
\toprule\noalign{}
\(k\) & \(R = 2k+1\) & \(N(k)\) & \(a_k := N(k) - N(k-1)\) \\
\midrule\noalign{}
\endhead
\bottomrule\noalign{}
\endlastfoot
0 & 1 & 1 & 1 \\
1 & 3 & 137 & 136 \\
2 & 5 & 1,545 & 1,408 \\
3 & 7 & 7,281 & 5,736 \\
4 & 9 & 22,409 & 15,128 \\
5 & 11 & 53,161 & 30,752 \\
6 & 13 & 108,081 & 54,920 \\
7 & 15 & 199,953 & 91,872 \\
8 & 17 & 337,417 & 137,464 \\
9 & 19 & 537,409 & 199,992 \\
10 & 21 & 818,145 & 280,736 \\
\end{longtable}
}

\subsubsection{§4.2 角の立方体}\label{ux89d2ux306eux7acbux65b9ux4f53}

\textbf{命題 4.1}。\emph{各 \(k \ge 0\) について、中心
\(c = (\pm k, \pm k, \pm k, \pm k)\) の16個の単位立方体は最遠頂点が球面
\(\partial B(R_k) = \partial B(2k+1)\)
上にある。これらを「角の立方体」と呼ぶ。}

\textbf{証明}。\(c = (\pm k)^4\) で全 \(i\) について
\(|c_i| = k\)、よって \(|c_i| + 1/2 = k + 1/2\)。
\[\sum_i (|c_i| + 1/2)^2 = 4 (k + 1/2)^2 = (2k + 1)^2 = R_k^2.\] 補題
3.1 の不等式は等号成立、最遠頂点は \(\partial B(R_k)\)
上に厳密に乗る。各象限に1個ずつ \(2^4 = 16\)
個の角の立方体が存在する。\(\square\)

\textbf{注意}。角の立方体の頂点が \(\partial B(R_k)\)
と厳密に接することは4次元で、かつ \(R = 2k+1\)
という選択で特別に成立する。次元 \(n\) の類似条件は
\[n \cdot (k + 1/2)^2 = R^2\] で、これが \(R = (2k+1) \sqrt{n}/2\)
で整数解を持つのは \(\sqrt{n}\) が有理数のとき、すなわち \(n\)
が完全平方数のときに限る。この性質を持つ最小の非自明次元は \(n = 4\)
で、このとき \(\sqrt{n} = 2\) かつ \(R = 2k+1\)。

\subsubsection{§4.3
漸近充填密度}\label{ux6f38ux8fd1ux5145ux586bux5bc6ux5ea6}

\textbf{充填密度}を
\[\rho(k) := \frac{N(k)}{V_4(R_k)} = \frac{N(k)}{(\pi^2/2) (2k+1)^4}\]
と定義する。これは球 \(B(R_k)\) のうち \(N(k)\)
個の単位立方体の和集合が占める体積分率。

\textbf{命題 4.2}。\emph{\(\lim_{k \to \infty} \rho(k) = 1\).}

\textbf{証明}。条件 \(\sum (|c_i| + 1/2)^2 \le R^2\) は領域
\(\Omega_R \subset \mathbb{R}^4\)
を定義する。標準的な凸体格子点定理により \(R \to \infty\) で
\[N(k) = \mathrm{Vol}(\Omega_R) + O(R^3).\] 直接計算（§5 で延期）により
\(\mathrm{Vol}(\Omega_R) = V_4(R) - O(R^3)\)。よって
\(N(k)/V_4(R) \to 1\)。\(\square\)

数値値の確認：

{\def\LTcaptype{none} % do not increment counter
\begin{longtable}[]{@{}rr@{}}
\toprule\noalign{}
\(k\) & \(\rho(k)\) \\
\midrule\noalign{}
\endhead
\bottomrule\noalign{}
\endlastfoot
0 & 0.2026 \\
1 & 0.3427 \\
2 & 0.5009 \\
5 & 0.7358 \\
10 & 0.8525 \\
20 & 0.9201 \\
30 & 0.9457 \\
40 & 0.9562 \\
50 & 0.9669 \\
60 & 0.9723 \\
\end{longtable}
}

収束は一様には速くない：\(k = 60\) で密度はまだ約
\(0.972\)。緩やかな収束は §5 で扱う境界補正により支配される。

\subsubsection{§4.4
体積不足とそのスケーリング}\label{ux4f53ux7a4dux4e0dux8db3ux3068ux305dux306eux30b9ux30b1ux30fcux30eaux30f3ux30b0}

\textbf{体積不足}を \[\Delta(R) := V_4(R) - N(k)\]
と定義する。「欠けた体積」を測る：連続4次元球体積と完全包含される単位立方体数の差。

本論文の主要結果（§5--§6 で証明）：
\[\Delta(R) = \frac{16\pi}{3} R^3 - 6\pi R^2 + O(R) \qquad (R \to \infty).\]

正規化形 \(c(k) := \Delta(R)/(2\pi^2 R^3)\) で主要定数：
\[\boxed{\;c := \lim_{k \to \infty} c(k) = \frac{16\pi/3}{2\pi^2} = \frac{8}{3\pi} \approx 0.84883.\;}\]

\(c(k)\) の数値値は下からこの極限に近づく：

{\def\LTcaptype{none} % do not increment counter
\begin{longtable}[]{@{}rr@{}}
\toprule\noalign{}
\(k\) & \(c(k)\) \\
\midrule\noalign{}
\endhead
\bottomrule\noalign{}
\endlastfoot
10 & 0.7745 \\
20 & 0.8192 \\
30 & 0.8276 \\
40 & 0.8311 \\
50 & 0.8350 \\
60 & 0.8380 \\
\end{longtable}
}

収束は緩やかだが、多項式フィットで \(k \ge 30\) で 0.024\% 以内（§5.4
参照）。

\subsubsection{§4.5
母関数（延期）}\label{ux6bcdux95a2ux6570ux5ef6ux671f}

Jacobi テータ関数による \(N(k)\) の母関数表現を §6 で展開する。§5
の漸近解析には不要。接続は別目的：\(N(k)\)
を平方和表現の古典的枠組み内に位置づけ、（§6.3
で延期される正確な符号補正係数を除き）厳密な閉形式を提供する。

\begin{center}\rule{0.5\linewidth}{0.5pt}\end{center}

\textbf{(§5--§7
は別途執筆、\href{draft_full_ja.md}{\texttt{draft\_full\_ja.md}}
に収録)}

\begin{center}\rule{0.5\linewidth}{0.5pt}\end{center}

\textbf{未確定事項（後で検討）}: - §5（漸近解析）と §6（Jacobi
接続）の整形済み本文（\texttt{draft\_full\_ja.md} 参照） - §1（序論）と
§7（結論）の本文

\begin{center}\rule{0.5\linewidth}{0.5pt}\end{center}

\subsection{\texorpdfstring{§5.
漸近解析：\(\Delta(R) \sim c\,R^3\)}{§5. 漸近解析：\textbackslash Delta(R) \textbackslash sim c\textbackslash,R\^{}3}}\label{ux6f38ux8fd1ux89e3ux6790deltar-sim-cr3}

\subsubsection{\texorpdfstring{§5.1 補助領域
\(\Omega_R\)}{§5.1 補助領域 \textbackslash Omega\_R}}\label{ux88dcux52a9ux9818ux57df-omega_r}

充填数 \(N(k)\) は領域
\[\Omega_R := \left\{x \in \mathbb{R}^4 : \sum_{i=1}^{4} (|x_i| + \tfrac{1}{2})^2 \le R^2\right\}\]
内の整数点を数える。標準的な凸体格子点定理により、有界領域
\(\Omega \subset \mathbb{R}^4\) に対し
\[\#(\mathbb{Z}^4 \cap \Omega) = \mathrm{Vol}(\Omega) + O(\partial\Omega)\]
が成り立つ。ここで \(O(\partial\Omega)\) は \(\Omega\)
の表面積に依存する誤差項。\(\Omega = \Omega_R\)（\(R = 2k+1\)）に対し、表面は
\(C^0\) 超曲面（\(x_i = 0\) で角を持つ）で、表面積は \(O(R^3)\)。よって
\[N(k) = \mathrm{Vol}(\Omega_R) + E(R), \qquad E(R) = O(R^3).\]

\(\Delta(R) = V_4(R) - N(k)\) の主要項は
\(V_4(R) - \mathrm{Vol}(\Omega_R)\)、格子点誤差 \(E(R)\)
は副主要項として §5.4 で扱う。

\subsubsection{\texorpdfstring{§5.2 包除原理による \(\Omega_R\)
の体積}{§5.2 包除原理による \textbackslash Omega\_R の体積}}\label{ux5305ux9664ux539fux7406ux306bux3088ux308b-omega_r-ux306eux4f53ux7a4d}

\(\Omega_R\) は \(x_i \mapsto -x_i\) の対称性で不変。よって正の象限
\(\{x : x_i \ge 0\}\) に制限し、結果を \(2^4 = 16\) 倍すれば良い：
\[\mathrm{Vol}(\Omega_R) = 16 \cdot \mathrm{Vol}(K_R), \qquad K_R := \{u \in \mathbb{R}^4 : u_i \ge 1/2,\ \textstyle\sum u_i^2 \le R^2\}.\]
ここで \(u_i := x_i + 1/2\)、\(u_i \ge 1/2\) は \(x_i \ge 0\) に対応。

\(K_R\) は球の正の象限 \(\{u : u_i \ge 0, \sum u_i^2 \le R^2\}\)（体積
\(V_4(R)/16\)）から、4本のスラブ \(\{u_i < 1/2\}\)
を除いた部分。包除原理：
\[\mathrm{Vol}(K_R) = \frac{V_4(R)}{16} - \sum_{i} V_i + \sum_{i<j} V_{ij} - \sum_{i<j<k} V_{ijk} + V_{1234}\]
ここで \(V_S\) は \(S\) 添字スラブの共通部分の体積。

\subsubsection{§5.3
スラブ体積の漸近評価}\label{ux30b9ux30e9ux30d6ux4f53ux7a4dux306eux6f38ux8fd1ux8a55ux4fa1}

\(R \gg 1\) で漸近展開：

\textbf{単一スラブ} \(V_i\)。\(u_i \in [0, 1/2]\)
で積分、残り3座標は正の象限 3-球：
\[V_i = \int_0^{1/2} \frac{1}{8} \cdot \frac{4\pi}{3} (R^2 - u_i^2)^{3/2}\, du_i = \frac{\pi R^3}{12} - \frac{\pi R}{96} + O(R^{-1}).\]

\textbf{ペアスラブ} \(V_{ij}\)：
\[V_{ij} = \int_0^{1/2}\int_0^{1/2} \frac{1}{4} \cdot \pi (R^2 - u_i^2 - u_j^2)\, du_i\, du_j = \frac{\pi R^2}{16} - \frac{\pi}{96} + O(R^{-2}).\]

\textbf{3重スラブ} \(V_{ijk}\)：
\[V_{ijk} = \int_{[0,1/2]^3} \frac{1}{2} \sqrt{R^2 - u_i^2 - u_j^2 - u_k^2}\, du = \frac{R}{16} + O(R^{-1}).\]

\textbf{4重スラブ} \(V_{1234} = 1/16\)。

\subsubsection{§5.4 体積公式}\label{ux4f53ux7a4dux516cux5f0f}

代入：
\[\mathrm{Vol}(K_R) = \frac{V_4(R)}{16} - \frac{\pi R^3}{3} + \frac{3\pi R^2}{8} - \frac{R}{4} + \frac{1}{16} + O(R^{-1}).\]

16 倍：
\[\boxed{\;\mathrm{Vol}(\Omega_R) = V_4(R) - \frac{16\pi R^3}{3} + 6\pi R^2 - 4 R + 1 + O(R^{-1}).\;}\]

したがって
\[V_4(R) - \mathrm{Vol}(\Omega_R) = \frac{16\pi R^3}{3} - 6\pi R^2 + 4R - 1 + O(R^{-1}).\]

\subsubsection{§5.5
格子点誤差と漸近係数}\label{ux683cux5b50ux70b9ux8aa4ux5deeux3068ux6f38ux8fd1ux4fc2ux6570}

\(N(k) = \mathrm{Vol}(\Omega_R) + E(R)\) と組み合わせて：
\[\Delta(R) = \frac{16\pi R^3}{3} - 6\pi R^2 + 4R - 1 + O(R^{-1}) - E(R).\]

格子点誤差 \(E(R) = O(R^3)\) は \(\Omega_R\)
に対し\textbf{振動的}であると期待される：凸体格子点問題の古典的結果により、\(|E(R)|\)
は表面積で抑えられるが、\(R\)
にわたる平均値ははるかに小さい（境界の正則性に応じて典型的に
\(O(R^{n-2+\alpha})\)、\(\alpha > 0\)）。我々の場合、\(\partial \Omega_R\)
は \(x_i = 0\) で角を持つ区分的滑らか曲面で、\(E(R)\) の主要 \(R^3\)
係数への寄与が平均ゼロになることを \textbf{予想}するが、証明は与えない。

この予想（§7 で数値的に検証）の下で、\(\Delta(R)\) の主要項は境界補正項
\(16\pi R^3/3\)： \[\Delta(R) \sim \frac{16\pi}{3}\, R^3.\]

正規化主要係数：
\[\boxed{\;c := \lim_{k \to \infty} \frac{\Delta(R)}{2\pi^2 R^3} = \frac{16\pi/3}{2\pi^2} = \frac{8}{3\pi} \approx 0.84883.\;}\]

幾何学的意味：\(\Delta(R)\) は \(\partial B(R)\) の表面積
\(S_3(R) = 2\pi^2 R^3\) と同じスケーリングを持ち、比例定数は
\(8/(3\pi) \approx 0.849\)。すなわち \(\Delta(R)\) は単位立方体幅の
\(\sim 8/(3\pi)\) 倍の厚みの境界層に集中する。

\begin{center}\rule{0.5\linewidth}{0.5pt}\end{center}

\subsection{§6. Lagrange--Jacobi
接続}\label{lagrangejacobi-ux63a5ux7d9a}

\subsubsection{§6.1
4個の正の奇数平方和への帰着}\label{ux500bux306eux6b63ux306eux5947ux6570ux5e73ux65b9ux548cux3078ux306eux5e30ux7740}

§3.2 の整数定式化から、充填不等式
\[\sum_i (|c_i| + 1/2)^2 \le R^2 \qquad (R = 2k+1, c \in \mathbb{Z}^4)\]
は \(p_i := 2|c_i| + 1\) により
\[p_1^2 + p_2^2 + p_3^2 + p_4^2 \le (4k+2)^2, \qquad p_i \in \{1, 3, 5, \ldots\}\]
となる。総数：
\[N(k) = \sum_{\substack{(p_i) \in \{1,3,5,\ldots\}^4 \\ \sum p_i^2 \le (4k+2)^2}} 2^{\#\{i : p_i > 1\}}.\]
\(2^{\#\{p_i > 1\}}\) は、\(p_i > 1\) なら \(c_i = \pm (p_i-1)/2\)
の2通り、\(p_i = 1\) なら \(c_i = 0\) の1通り、という符号自由度を表す。

\subsubsection{§6.2 Jacobi
の四平方定理}\label{jacobi-ux306eux56dbux5e73ux65b9ux5b9aux7406}

古典的 Jacobi 四平方定理：\(N\) を順序付き4個の整数平方の和（0
と負も許容）として表す方法の数 \(r_4(N)\) は
\[r_4(N) = \begin{cases} 8\,\sigma(N) & 4 \nmid N \\ 24\,\sigma(N_{\mathrm{odd}}) & 4 \mid N \end{cases}\]
ここで \(\sigma(N)\) は約数和、\(N_{\mathrm{odd}}\) は \(N\)
の最大奇成分。

我々の目的に必要なのは、4個の\textbf{奇数}整数平方（正負）の和としての表現。\(p_i\)
が全て奇数なら \(p_i^2 \equiv 1 \pmod 8\) より
\(\sum p_i^2 \equiv 4 \pmod 8\)。逆に \(N \equiv 4 \pmod 8\) なら、\(N\)
を4平方和として表す任意の表現は全平方が奇数（\(\bmod 8\)
解析より）。よって
\[\#\{(p_i) \in \mathbb{Z}^4 \text{ 全奇} : \sum p_i^2 = N\} = \begin{cases} r_4(N) & N \equiv 4 \pmod 8 \\ 0 & \text{その他} \end{cases}\]

\subsubsection{§6.3
閉形式表現と構造的恒等式}\label{ux9589ux5f62ux5f0fux8868ux73feux3068ux69cbux9020ux7684ux6052ux7b49ux5f0f}

\(N \le (4k+2)^2\), \(N \equiv 4 \pmod 8\)
で総和すると、符号によらない全奇整数4組の総数を得る。これを正の奇数で重み
\(2^{\#\{p_i > 1\}}\)
の数え上げに変換するには、より丁寧な包除原理が必要である。\textbf{後で検討}：\(N(k)\)
の \(\sum_N r_4(N)\) による正確な閉形式表現と、\(p_i = 1\)
部分集合に対応する境界補正項の決定。

正確に述べられるのは以下の構造的恒等式：
\[T(M) := \sum_{\substack{N \le M \\ N \equiv 4 \pmod 8}} r_4(N) = \sum_{\substack{N \le M \\ N \equiv 4 \pmod 8}} 24\,\sigma(N_{\mathrm{odd}}).\]
\(T(M)\) は全 \((p_i) \in \mathbb{Z}^4\)（全奇）で \(\sum p_i^2 \le M\)
を満たすものを数える。\(T((4k+2)^2)\) と \(N(k)\) の関係は
\(\{p_i = 1\}\) 事象に対する包除原理を含み、より小さい引数での
\(T\)-値で完全に決定される。数値的対応（§7 で確認）は、\(N(k)\) が
\(T\)-値による正確な閉形式表現を持つことを支持する。正確な係数の記号的決定は今後の課題。

\subsubsection{§6.4
4次元の特殊性}\label{ux6b21ux5143ux306eux7279ux6b8aux6027}

閉形式
\(r_4(N) = 8\sigma(N)\)（\(4 \nmid N\)）は4次元の特徴的な現象である。\(n \neq 1, 2, 4, 8\)
では、表現数 \(r_n(N)\) は単純な閉形式を持たない：\(r_3(N)\) は
Hurwitz--Kronecker 類数を含み、\(r_5(N), r_6(N), \ldots\)
は普遍的閉形式を持たないモジュラー形式を含む。4次元の単位立方体充填問題がこの構造的単純性を継承するという事実が、§5
の精密な漸近係数 \(c = 8/(3\pi)\) の決定の理由である。

\begin{center}\rule{0.5\linewidth}{0.5pt}\end{center}

\subsection{§7. 数値検証}\label{ux6570ux5024ux691cux8a3c}

\subsubsection{§7.1 計算結果}\label{ux8a08ux7b97ux7d50ux679c}

§3.4 の対称性削減列挙により、\(N(k)\) を \(k = 0, 1, \ldots, 60\)
まで厳密計算した。完全データは
\href{../../computations/results/N_table_k0_60.tsv}{\texttt{computations/results/N\_table\_k0\_60.tsv}}
を参照。代表値：

{\def\LTcaptype{none} % do not increment counter
\begin{longtable}[]{@{}rrrrrr@{}}
\toprule\noalign{}
\(k\) & \(R\) & \(N(k)\) & \(V_4(R)\) & \(\Delta(R)\) & \(c(k)\) \\
\midrule\noalign{}
\endhead
\bottomrule\noalign{}
\endlastfoot
10 & 21 & 818,145 & 959,725.27 & 141,580.27 & 0.7745 \\
20 & 41 & 12,830,145 & 13,944,571.60 & 1,114,426.60 & 0.8192 \\
30 & 61 & 64,618,369 & 68,326,486.64 & 3,708,117.64 & 0.8276 \\
40 & 81 & 197,955,121 & 207,019,330.71 & 9,064,209.71 & 0.8311 \\
50 & 101 & 496,535,873 & 513,517,495.84 & 16,981,622.84 & 0.8350 \\
60 & 121 & 1,028,515,513 & 1,057,818,677.67 & 29,303,164.67 & 0.8380 \\
\end{longtable}
}

充填密度 \(\rho(k) = N(k)/V_4(R)\) は命題 4.2 の通り下方から \(1\)
に漸近。

\subsubsection{§7.2
多項式フィット}\label{ux591aux9805ux5f0fux30d5ux30a3ux30c3ux30c8}

\(\Delta(R) = c_3 R^3 + c_2 R^2 + c_1 R + c_0\) を
\(k_{\min} \le k \le 60\) にフィット：

{\def\LTcaptype{none} % do not increment counter
\begin{longtable}[]{@{}
  >{\raggedleft\arraybackslash}p{(\linewidth - 10\tabcolsep) * \real{0.1667}}
  >{\raggedleft\arraybackslash}p{(\linewidth - 10\tabcolsep) * \real{0.1667}}
  >{\raggedleft\arraybackslash}p{(\linewidth - 10\tabcolsep) * \real{0.1667}}
  >{\raggedleft\arraybackslash}p{(\linewidth - 10\tabcolsep) * \real{0.1667}}
  >{\raggedleft\arraybackslash}p{(\linewidth - 10\tabcolsep) * \real{0.1667}}
  >{\raggedleft\arraybackslash}p{(\linewidth - 10\tabcolsep) * \real{0.1667}}@{}}
\toprule\noalign{}
\begin{minipage}[b]{\linewidth}\raggedleft
\(k_{\min}\)
\end{minipage} & \begin{minipage}[b]{\linewidth}\raggedleft
\(c_3\)
\end{minipage} & \begin{minipage}[b]{\linewidth}\raggedleft
\(c_3/(2\pi^2)\)
\end{minipage} & \begin{minipage}[b]{\linewidth}\raggedleft
\(8/(3\pi)\) との差
\end{minipage} & \begin{minipage}[b]{\linewidth}\raggedleft
\(c_2\)
\end{minipage} & \begin{minipage}[b]{\linewidth}\raggedleft
\(c_2/(-6\pi)\)
\end{minipage} \\
\midrule\noalign{}
\endhead
\bottomrule\noalign{}
\endlastfoot
10 & 16.7987 & 0.85103 & \(+0.0022\) & \(-37.36\) & \(1.98\) \\
20 & 16.8140 & 0.85181 & \(+0.0030\) & \(-40.82\) & \(2.16\) \\
30 & 16.7511 & 0.84862 & \(-0.00021\) & \(-22.22\) & \(1.18\) \\
\end{longtable}
}

主要係数 \(c_3\) は \(k_{\min} = 30\) で解析値 \(16\pi/3 = 16.7552\) と
0.024\% 以内で一致。副主要係数 \(c_2\) は \(-6\pi = -18.85\)
とオーダー一致するが、解析的に解決していない \(E(R)\) の \(R^2\)
寄与に敏感で大きく揺らぐ。\textbf{後で検討}：\(E(R)\)
のより精密な漸近解析。

\subsubsection{§7.3 結論}\label{ux7d50ux8ad6}

数値証拠は主要漸近係数の解析的決定を支持。格子点誤差 \(E(R)\)
は計算精度の範囲内でゼロ平均を持ち、§5.5 の予想を正当化する。

\begin{center}\rule{0.5\linewidth}{0.5pt}\end{center}

\subsection{§8. 結論}\label{ux7d50ux8ad6-1}

以下を確立した：

\begin{enumerate}
\def\labelenumi{\arabic{enumi}.}
\tightlist
\item
  \textbf{正確な定式化}：整数格子点中心の単位立方体が \(B(R)\)
  に包含される条件は \(\sum_i (|c_i|+1/2)^2 \le R^2\)。\(R = 2k+1\)
  に対する \(N(k)\) は対称性削減列挙で正確に計算可能。
\item
  \textbf{数値}：\(N(k)\) を \(k = 0, 1, \ldots, 60\)
  まで厳密計算、\(N(60) = 1{,}028{,}515{,}513\)。
\item
  \textbf{漸近係数}：体積不足は
  \[\Delta(R) = \frac{16\pi}{3}\, R^3 - 6\pi\, R^2 + 4R - 1 + O(R^{-1})\]
  を満たす。\(R^3\) および \(R^2\) 係数は補助領域 \(\Omega_R\)
  への包除原理から得られ、低次項は体積展開と格子点誤差 \(E(R)\)
  の両方を含む。
\item
  \textbf{正規化定数}：\(c := \lim_{k \to \infty} \Delta(R)/(2\pi^2 R^3) = 8/(3\pi) \approx 0.84883\)。
\item
  \textbf{数論的構造}：充填不等式は4個の正奇数平方和への制約と同値で、Lagrange--Jacobi
  枠組み内の問題。\(N(k)\) の \(r_4(N)\) による正確な閉形式表現は
  \(\{p_i = 1\}\) 事象に対する包除原理を含み、正確な係数は今後の課題。
\item
  \textbf{数値検証}：\(k \ge 30\)
  の多項式最小二乗フィットで主要係数が解析値と 0.024\% 以内で一致。
\end{enumerate}

4次元は閉形式 Jacobi \(r_4\)
公式により特殊な構造的単純性を享受する。高次元 \(n = 5, 6, 7, \ldots\)
では類似の充填問題が定式化可能だが、漸近係数は単純な閉形式を持たないモジュラー形式の係数を含む。

主要な未解決問題： - \(N(k)\) の \(r_4(N)\)
による閉形式表現の精密な記号的決定（境界補正の正確な係数を含む） -
格子点誤差 \(E(R)\) のより精密な漸近解析 - 次元 \(n \neq 4\)
および非立方充填セルへの拡張

\begin{center}\rule{0.5\linewidth}{0.5pt}\end{center}

\subsection{\texorpdfstring{参考文献（選定、\textbf{後で検討}）}{参考文献（選定、後で検討）}}\label{ux53c2ux8003ux6587ux732eux9078ux5b9aux5f8cux3067ux691cux8a0e}

\begin{itemize}
\tightlist
\item
  C. F. Gauss, \emph{Disquisitiones Arithmeticae} (1801).
\item
  J.-L. Lagrange, ``Démonstration d'un théorème d'arithmétique,''
  \emph{Nouv. Mém. Acad. Berlin} (1770).
\item
  C. G. J. Jacobi, \emph{Fundamenta Nova Theoriae Functionum
  Ellipticarum} (1829).
\item
  K. Mahler, \emph{Lectures on Diophantine Approximations} (1961).
\item
  G. H. Hardy and E. M. Wright, \emph{An Introduction to the Theory of
  Numbers}, 第6版 (Oxford, 2008).
\item
  M. N. Huxley, \emph{Area, Lattice Points, and Exponential Sums}
  (Oxford, 1996).
\end{itemize}

\end{document}
